% =============================================================================
% Chapter 2: Plain-English Ontology
% Input-ready fragment for the lean toy-model research proposal.
% This file intentionally contains no preamble or document environment.
% =============================================================================

\chapter{Plain-English Ontology}
\label{ch:plain-english-ontology}

Ontology describes what the candidate model says exists and how its familiar
phenomena are to be interpreted. It does not establish that the required
solutions exist. Some elements below are defining postulates, some are physical
interpretations of quantities that a brane observer would measure, and some are
open hypotheses that later calculations may confirm or reject. The canonical
ontology and closure ledger remains the authority for exact definitions,
branch conditions, symbols, and failure criteria; this chapter supplies the
readable mental picture needed to understand the research program.

The proposal is deliberately a one-medium construction. Light, particles,
charge, gravity, magnetism, radiation, inertia, and return are not assigned to
independent substances that can be tuned separately. They are intended to be
different behaviors of one material system. The central scientific burden is
therefore not merely to make each interpretation plausible in isolation, but to
show that one frozen parent medium can realize all of them together.

\section{Arena and one fundamental medium}

The parent arena has one preferred time coordinate and four spatial coordinates.
A brane observer labels an ordinary position by

\[
\mathbf x=(x,y,z),
\]

while a complete parent-space position is

\[
\mathbf X=(\mathbf x,w).
\]

The coordinate \(w\) is normal to the brane. The present candidate takes the
four-dimensional spatial geometry to be flat and Euclidean, with preferred
Newtonian time. This arena is a postulate of the toy model rather than a result
to be derived. Changing it later would alter the meaning of distance,
reflection, wave propagation, Green functions, and throat contact, so it would
constitute a new architectural branch rather than a harmless change of
notation.

A brane-bound observer normally records only \(\mathbf x\). Two structures can
therefore have the same projected three-dimensional position while remaining
separated in \(w\). That distinction becomes physically important for finite
throats: projected overlap of their apparent mouths does not by itself imply
that their full four-dimensional cores touch.

The brane is presently assumed to be globally two-sided and co-orientable. In
plain language, a continuous normal direction can be chosen throughout the
brane, allowing the two sides to be compared consistently. Calling them \(+w\)
and \(-w\) is conventional, but the availability of two distinguishable sides
is part of the candidate geometry. The \emph{sign} of the normal is used by the
charge ontology, while the unsigned normal axis is enough to organize ordinary
light guidance.

Everything in this arena is made from one compressible condensate. At the
mean-field level it is represented schematically by

\[
\psi=\sqrt n\,e^{i\theta},
\]

where \(n\) is constituent density and the phase \(\theta\) can generate
material flow on a conservative branch. This field is a description of the
substrate, not a claim that observed particles are collections of conventional
massive constituents whose rest masses add up to the observed particle mass.
The observed particle, if the model succeeds, is an emergent configuration of
the medium.

Total constituent content obeys a parent conservation law,

\[
\partial_t n+\nabla_4\!\cdot\mathbf J_n=0.
\]

Material may move, compress, change its degree of order, and exchange energy
among resolved and unresolved modes, but it may not be created from an empty
bulk or disappear into an external supply. An apparent source or sink seen by a
three-dimensional observer must ultimately be a projection of four-dimensional
transport or a conversion between states of the same medium.

Because the arena includes a preferred material frame, Lorentz symmetry is not
fundamental. It could emerge approximately only if the relevant long-wavelength
branches acquire sufficiently aligned speeds and couplings. The speeds of
transverse light, longitudinal brane motion, bulk sound, electric response, and
gravity-related response are therefore distinct until the common-medium
calculation proves an acceptable relation among them.

\section{Ordered brane and de-structured bulk}

An order variable \(\chi_B\) distinguishes two material states. Where
\(\chi_B\simeq1\), the medium is ordered, brane-like, and able to carry the
resolved shear structure required by the model. Where \(\chi_B\simeq0\), the
same material is de-structured and bulk-like. The brane and bulk are therefore
not two substances. They are two states of one globally conserved substance,
and their exchange routes are not symmetric. Ordered material may de-structure
as a localized one-way throat drain carries it into the bulk, but it does not
return through the throat. Because the finite-thickness brane is porous, bulk
material can instead reorganize into the ordered brane state across the brane;
for a projected observer, the positive bulk-to-brane part of that distributed
exchange is encoded in \(S_{\mathrm{leak}}\).

Observed space is identified with a finite ordered slab centered on a local
mid-surface. It is not an infinitesimally thin mathematical sheet. Bulk-like
material lies on both sides, and the slab has two interfaces separated by a
characteristic thickness \(H_{\mathrm{br}}\). The local mid-surface is
relational: translating the entire slab through a homogeneous parent medium
changes the coordinate origin but not the physical state. What matters is
displacement relative to the slab's own center and interfaces.

While the two interfaces remain single-valued surfaces over the brane, their
displacements can be written as \(h_+\) and \(h_-\). Two combinations are useful:

\[
h_m=\frac{h_++h_-}{2},
\qquad
h_t=\frac{h_+-h_-}{2}.
\]

The first is the leading mid-surface displacement. It reverses under reflection
through the slab and is therefore reflection odd. The second changes the local
half-thickness and is reflection even. These parity labels organize the
far-field calculation, but they do not guarantee that either variable is a
complete physical eigenmode. Density, order, phase, longitudinal motion, bulk
response, and reservoir variables may mix with both.

The two-interface description is also limited to the graph regime. It is useful
for the stable background, weak waves, far fields, and separated deformations.
It cannot represent a throat after the ordered region develops a tube, overhang,
neck, fold, reconnection, pinch-off, or other topology change. In that regime
the full parent fields and localized core diagnostics must replace the interface
coordinates.

The ordered slab is intended to support in-plane shear. That property is what
makes transverse light possible. The de-structured bulk is not assumed to carry
the same resolved shear branch, so the contrast between the two states may allow
the slab to act as a waveguide. Finite thickness also supplies a material length
scale that can affect dispersion, interface mixing, electric response, throat
size, and the crossover from long-range behavior to nonlinear core physics.
It is therefore a load-bearing part of the ontology rather than a decorative
embedding picture.

The exact dynamics of order remain a major branch choice. In a reversible
realization, changes in \(\chi_B\) behave as inertial material dynamics and carry
energy and momentum coherently. In a relaxational realization, loss and recovery
of order may produce entropy and require explicit internal heat or reservoir
variables. A mixed realization may contain both reversible storage and damping
into identified internal degrees of freedom. What is not permitted is to use a
conservative energy whenever it is convenient and then add untracked damping
whenever stability or transport becomes difficult. The selected branch must
close its own material, momentum, energy, and, where applicable, entropy ledger.

The bulk consequently functions as an internal reservoir. It may store
compression, flow, internal energy, disorder, heat, or other unresolved
excitations of the same substance. It is neither an empty vacuum that pulls on
the brane nor an inexhaustible battery. Any sustained throat process must alter
that reservoir and eventually be solved together with the local brane and
throat configuration.

\section{Light in the brane}

Light is proposed to be transverse shear motion guided by the ordered slab. For
a wave traveling along one brane direction, its displacement lies within the
brane but perpendicular to the direction of travel. In the four-dimensional
parent space, excluding the one normal direction and the one propagation
direction leaves two independent tangential polarization directions. Under the
currently supplied homogeneous, isotropic quadratic brane action, that
geometrical count is realized as two transverse modes with characteristic speed

\[
c_\gamma^2=\frac{\mu_R}{\rho_{\mathrm{br}}},
\]

where \(\mu_R\) is an effective shear modulus and \(\rho_{\mathrm{br}}\) is an
effective brane inertia density. The same quadratic system contains one separate
in-plane longitudinal mode and no transverse--longitudinal cross-block at that
order. This is a conditional result of the stated homogeneous action, not yet a
proof of confinement on a finite nonuniform slab.

Three related objects must remain distinct. A \emph{background light wave} is a
small-amplitude transverse excitation of the throat-free brane; ordinary plane
waves and spreading linear packets belong here. A \emph{free photon candidate}
is a much stronger object: a finite-energy traveling packet localized in all
three brane directions, guided across \(w\), and stable or sufficiently
long-lived. A \emph{throat support mode} is a transverse standing mode bound to
or resonant around a particle throat. It may help hold the throat open even if
free light does not form a self-bound packet in empty brane space.

Guidance across \(w\) and localization within the brane are separate problems.
The finite slab may prevent a transverse mode from leaking normally into the
bulk while its envelope still spreads in \(x\), \(y\), or along its propagation
direction. The stronger photon hypothesis therefore asks whether finite-thickness
dispersion and nonlinear material response can maintain a three-dimensional
packet.

The intended packet remains predominantly transverse. Its intensity may excite
a reflection-even dressing made from some combination of longitudinal
compression, thickness change, density change, and order change. If that helper
response is determined by the transverse intensity rather than independently
selectable, it can alter the local propagation environment without becoming a
third photon polarization. Because the dressing is reflection even, ordinary
free light can use the slab's normal structure without acquiring the leading
reflection-odd electric monopole assigned to charged throats.

The longitudinal branch is therefore neither charge nor a mode to be discarded
merely because exact Maxwell theory would not contain it in the same form. It
may contribute to the slaved helper response, remain as an independent guided
wave, leak into bulk sound, or become an unwanted radiation and drag channel.
The complete finite-slab spectrum must decide. A viable light sector must also
control polarization splitting, bulk leakage, defect-induced conversion,
dispersion, and any nonlinear wake. Failure of the self-bound packet would
leave a possible classical two-polarization light-wave sector, but it would not
complete the intended ontology of an individual freely propagating photon.

\section{Particles as finite supported throats}

The ontology contains no fundamental point particles. A particle is proposed to
be a finite nonlinear throat, puncture, or defect in which the ordered brane
extends into the bulk. A candidate branch may be labeled

\[
\mathcal T_{\alpha,s},
\qquad
s\in\{+1,-1\}.
\]

The label \(\alpha\) contains orientation-independent species information: the
support-mode family, core topology, boundary data, equilibrium size and shape,
transport and conversion pattern, and any additional stable or metastable
internal structure. The sign \(s\) records whether the throat is oriented toward
\(+w\) or \(-w\). Reversing \(s\) does not by itself turn one species into an
unrelated species.

The complete particle is more than its apparent mouth. It includes the full
four-dimensional throat geometry; the internal transverse standing mode; the
ordered rim, collar, sheath, cavity wall, or other region in which support stress
is concentrated; surrounding flow, pressure, density, order, and brane
deformation; the conversion region; its long-range fields; and the return and
reservoir response required by the selected dynamical branch. Here return is the
separate distributed porous-brane \(S_{\mathrm{leak}}\) process surrounding the
object, not reverse flow through its throat. A brane observer may assign the
object one projected location, but the physical core occupies a finite region of
\((\mathbf x,w)\).

A trapped transverse standing mode is proposed to help hold the aperture open.
Its stress is structural. The model does not identify the support-mode energy
with observed gravitational or inertial mass, and it does not assume a relation
such as \(E_{\mathrm{support}}=Mc^2\) as an input. The equilibrium throat must
instead arise from a balance among support stress, brane tension, material
transport, order conversion, bulk backpressure, and the reservoir state.

A linear bound mode on a provisional throat geometry is only a spectral seed.
The actual particle must be a self-consistent nonlinear solution in which the
oscillating support field and the time-averaged geometry, flow, conversion, and
reservoir loading agree. The throat can therefore be stationary in its average
shape and rates while its support field remains periodic. Its stability may
require a Floquet or retarded-response analysis rather than the Hessian of a
strictly static energy.

Material throughput and order conversion are related but distinct parts of the
throat. The net constituent flow across an aperture, denoted schematically by
\(Q_n\) on a selected outward-drain branch, measures transport. A gross local
ordered-to-bulk conversion diagnostic, denoted \(Q_\chi\), measures loss of
resolved brane order in the throat region. They need not be equal, and their
precise definitions depend on the chosen reversible, relaxational, or mixed
order branch. Neither quantity is a return rate: material that passes through
the throat does not come back through that aperture. Return is instead the
separate distributed reordering of bulk material into the porous brane,
represented by the positive bulk-to-brane contribution to
\(S_{\mathrm{leak}}\).

Different particle species may correspond to different solved throat branches,
not merely to different electric orientations. Species structure could reside
in distinct support modes, knot or core topology, composite geometry, or other
orientation-independent data. Consequently, reversing an electron-like throat
produces the leading positron-like candidate on the same internal branch; it
does not produce a proton-like object. A proton-like object would require a
different or composite branch.

The particle hypothesis remains open. No solution has yet jointly established a
finite supported geometry, localized or sufficiently long-lived support mode,
acceptable throughput and conversion, stable source profiles, reservoir loading,
and perturbative or Floquet stability. Later throat work is intended to decide
whether this picture can be realized by the frozen parent medium.

\section{Charge, gravity, magnetism, and radiation}

\subsection{Electric orientation and finite-thickness charge}

Electric sign is identified with throat orientation:

\[
s=+1\quad\text{for a throat toward }+w,
\qquad
s=-1\quad\text{for a throat toward }-w.
\]

The sign is an intrinsic property of an oriented defect, but it does not by
itself determine the magnitude of the long-range coupling, the particle species,
or the complete antiparticle operation. For a finite slab, an oriented throat
generally perturbs both interfaces. Its leading mid-surface-like disturbance is
reflection odd and reverses with \(s\), while a thickness-like disturbance is
reflection even and can remain the same for both orientations.

The long-range electric carrier is therefore not assumed to be the bare height
of one interface. It must emerge as a healthy odd eigencombination of every
relevant coupled variable. To produce a Coulomb-like potential in three brane
dimensions, the static long-wavelength operator must contain a branch with the
appropriate leading \(k^2\) stiffness. Reflection oddness or gaplessness alone
is insufficient; a different low-wave-number stiffness would produce a
different range law.

The intended force is a material compatibility effect. A throat imposes an
oriented boundary condition on the finite slab. Same-oriented defects should
compound the departure from the preferred brane state, while opposite-oriented
defects should partly cancel it. The actual force sign must be derived from the
correct equilibrium boundary ensemble or, for periodic or driven throats, from
cycle-averaged stress, momentum flux, conversion, and reservoir response. It
cannot be fixed merely by naming a variable ``electric.''

At long range the desired organization is

\[
q_a=s_a g_a,
\qquad g_a\ge0,
\]

where \(s_a\) carries all electric sign and \(g_a\) is a derived coupling
magnitude. Approximate unit-charge universality requires ordinary one-charge
throat branches to produce essentially the same leading \(g_a\), even when their
internal species structures differ. A common slab thickness helps but does not
prove this; the core boundary data must also be universal, quantized,
topologically controlled, or sufficiently decoupled from the far field.

A binary orientation label also does not automatically provide additive charge
conservation. The model must derive a conserved oriented count or current and
show how it appears in a projected Gauss-like law. Likewise, the complete
antiparticle candidate may require more than geometric reflection. In general
the target operation is

\[
\mathcal C_{\mathrm{th}}
=
\mathcal I_{\mathrm{internal}}\circ\mathcal R_w,
\]

where \(\mathcal R_w\) reverses the embedding orientation and
\(\mathcal I_{\mathrm{internal}}\) transforms any additional odd phase,
circulation, chirality, support-mode convention, topological datum, or conversion
current. Whether the internal transformation is trivial is an output of each
solved species branch.

Finite thickness prevents electrical cancellation from being mistaken for
material disappearance. Opposite orientations may occupy the same projected
position while their cores remain on opposite sides of the slab. Even if the
complete odd electric source cancels, the even thickness deformation and the
internal support structures may remain. The required logical separation is

\[
\text{odd electric cancellation}
\not\Rightarrow
\text{even material cancellation}
\not\Rightarrow
\text{core contact}
\not\Rightarrow
\text{annihilation}.
\]

A same-species throat--antithroat pair may eventually annihilate, reconnect,
scatter, or form a temporary neutral state. Oppositely charged objects of
different species may instead retain unmatched internal structure and form a
stable neutral composite. Electric neutrality alone does not determine the
short-range outcome.

\subsection{Gravity as a source-side material response}

Gravity is proposed to arise from the complete material configuration associated
with an open throat, not from the trapped standing-wave energy alone. The phrase
``gravity is inflow'' is a useful first picture for the localized one-way throat
drain, but the relevant source may also include pressure, acceleration of the
medium, ordered stress, de-structuring, the separate distributed porous-brane
\(S_{\mathrm{leak}}\) return, bulk compression, interface response, and
reservoir backreaction.

The environmental medium must first produce a probe-independent source field.
One fixed operational calibration then maps that field into the acceleration
units used by a brane observer. Only after that calibration is frozen is the
passive response of an arbitrary test throat calculated. A reference object may
set units once, but the calibration may not absorb species-dependent response
and make universal free fall true by definition. The calibrated gravity
observable is therefore not assumed to be numerically identical to the raw
four-dimensional material velocity.

Electric orientation and gravitational sign obey different tests. Under the
complete reflected throat operation and reflected environmental data, the
ordinary gravity source should remain the same when \(s\) reverses. Gravity is
targeted to be orientation even, while charge is orientation odd. That parity
test does not establish attraction: the local radial sign must separately be
shown to point inward on the ordinary branch.

The separate distributed return---bulk material reorganizing into the ordered
state across the porous brane, represented by the positive bulk-to-brane
contribution to \(S_{\mathrm{leak}}\)---can change the enclosed gravity source
with distance. It may
leave an intermediate inverse-square regime, screen the source beyond a return
scale, overshoot, or even reverse the effective sign. These outcomes must remain
visible rather than being removed by an absolute-value convention. The ontology
therefore distinguishes a local active source from the scale-dependent effective
source after return is included. Earlier effective calculations coupling a
localized sink to a separate distributed porous-brane return support
inverse-square scaling and leading attraction under stated assumptions, but the
full source reduction, calibration, parity closure, dynamics, and return profile
remain open.

\subsection{Magnetism and radiation from motion}

Electromagnetic magnetism is proposed to be the sign-dependent transverse
response of the same oriented throat when it moves. At low speed the source has
the schematic form

\[
\mathbf J_T\propto s\mathbf V.
\]

Reversing the throat orientation or the velocity reverses the magnetic source.
The coefficient, sign, tensor structure, and retarded response must ultimately
come from the solved moving throat rather than from an independently assigned
magnetic rule.

The intended electromagnetic hierarchy is

\[
O(v^0):\ \text{static electricity},
\qquad
O(v):\ \text{magnetism},
\qquad
O(a):\ \text{transverse radiation}.
\]

One conserved oriented current and one shared coupling structure must control
all three orders. A static electric coefficient, a separate magnetic
normalization, and a third radiation coupling would not establish one
medium-based electromagnetic sector.

Several motion-induced effects must remain separate. Electromagnetic magnetism
is sign dependent. The inertial wake or rebuilding of the surrounding medium is
primarily sign independent. Motion may also alter throat throughput and thereby
modify gravity. Gravitomagnetism is vorticity or frame-dragging-like structure
in the gravity-associated flow and is not the same phenomenon as electromagnetic
magnetism.

Acceleration must be tested for transfer of energy and momentum into both
transverse light polarizations. The same event may also emit into the odd
electric branch, even thickness modes, the longitudinal brane mode, bulk sound,
or order and reservoir excitations. Those additional channels are calculable
predictions. They may not be omitted simply because they have no counterpart in
exact Maxwell theory.

\section{Inertia and the complete dressed particle}

Inertia is proposed to be the resistance of the complete dressed throat to
acceleration. The particle is not a rigid mouth moving through an otherwise
unchanged background. It carries an oscillating support mode, a deformed brane,
near-field flow, pressure and density disturbances, a conversion region, and a
reservoir response. Changing the throat's velocity requires some or all of that
configuration to be rebuilt.

The trapped standing mode may need to retime its counterpropagating components,
shift phase and nodes, or change its resonant geometry as the throat accelerates.
Momentum must be transferred into that reconfiguration even if the local
transverse wave speed remains unchanged. The surrounding medium must also adjust
its flow, pressure, order, and deformation. Reversible rebuilding can contribute
an added-inertia-like response; outgoing waves and irreversible lag instead
appear as radiation or damping and must be accounted for separately.

The word \emph{mass} therefore refers to three distinct observational roles:

\begin{center}
\begin{tabularx}{\textwidth}{@{}L{0.24\textwidth}X@{}}
\toprule
\textbf{Role} & \textbf{Meaning in this proposal} \\
\midrule
Active gravitational mass & The source strength inferred from the local gravity field after the source-side reduction and one fixed calibration. \\
Passive gravitational mass & The low-frequency coefficient governing how a test throat responds to an already defined external gravity field. \\
Inertial mass & The low-frequency coefficient governing resistance of the complete dressed throat to acceleration. \\
\bottomrule
\end{tabularx}
\end{center}

These roles may arise from different pieces of the same solved object. Their
equality is not assumed. The ordinary branch must first produce positive active,
passive, and inertial behavior under fixed conventions. It must then show that
the appropriate low-energy ratios are sufficiently universal across species.
Universality of a wrong-sign, singular, or unstable response would not count as
success.

The calculation depends on the dynamical branch. In a closed reversible theory,
inertia may be extracted from the velocity dependence of the total momentum of
the support mode, geometry, near fields, and conversion sector. In a
relaxational or mixed theory, the response is frequency dependent and must
include damping, memory, internal heat or reservoir variables, and entropy
production. An unexcited reservoir may not deliver net work during a small
periodic perturbation; the retarded response must be passive.

The preferred material frame also creates a possible failure channel. Uniform
motion may emit into any branch whose phase velocity permits it, and a
relaxational medium may produce friction even when no propagating wave is
radiated. The full spectrum and low-frequency response must therefore be checked
for Cherenkov-like thresholds, modes with vanishing critical velocity,
conversion lag, and other forms of vacuum drag. The model has identified
plausible sources of inertia, but no moving supported-throat family has yet
derived a positive inertial response, acceptable drag, or the required
universality.

\section{Global conservation, return, and reservoir closure}

A throat provides a localized, one-way drain from the ordered brane. Material
passing through the throat mouth de-structures into the bulk state; it does not
return through the throat. Net aperture throughput tracks where material moves,
while gross local order conversion tracks how much resolved brane order is lost
in the throat region. These are not automatically the same rate, and neither is
the distributed return rate. A correct branch must derive all three without
hiding them inside one adjustable ``drain strength.''

Return is a separate, spatially distributed process. The finite-thickness brane
is porous, so bulk material can cross into the resolved brane region and
reorganize into the ordered brane state away from the throat mouth. For a
projected brane observer, constituent continuity has the form

\[
\partial_t\rho_{\mathrm{brane}}
+\nabla_3\!\cdot\!\mathbf j_{\mathrm{brane}}
=S_{\mathrm{leak}}.
\]

The positive bulk-to-brane contribution to the signed exchange term
\(S_{\mathrm{leak}}\) represents distributed return. The full
\(S_{\mathrm{leak}}\) is not synonymous with return: it is the net transverse
exchange seen by the projected observer and may have either sign. In an explicit
two-state order ledger, transport and order conversion are refined further, but
none of these descriptions permits a local down-and-back-up throat loop.

In a closed realization the material cycle is

\[
\begin{aligned}
\text{ordered brane material}
&\longrightarrow \text{localized one-way throat drainage}\\
&\longrightarrow \text{throat-region de-structuring and bulk reservoir}\\
&\longrightarrow \text{distributed }S_{\mathrm{leak}}
  \text{ reordering across the porous brane}\\
&\longrightarrow \text{ordered brane material}.
\end{aligned}
\]

The bulk may store density and pressure changes, flow, internal energy, disorder,
heat, or other excitations. The return leg is not reverse throat flow.
Backpressure or chemical-potential imbalance within the same medium may drive
bulk material to reorganize into the ordered state across the porous brane. It
may not be supplied by an unchanged external pump.

The local throat drain and the global porous-brane return--reservoir system are
nevertheless one feedback problem. A provisional local solution produces
one-way throughput, conversion, stress, emitted waves, and energy transfer.
Those outputs alter the global bulk, distributed return profile, and reservoir
state. The changed reservoir determines the distributed \(S_{\mathrm{leak}}\)
return across the brane and also alters the backpressure and asymptotic data seen
by the throat; it does not send material back through the throat. Closure
requires iteration until the local throat and global solution reproduce the same
boundary conditions, fluxes, distributed return profile, and reservoir state. A
local throat is not fully solved if its own long-term operation would invalidate
the environmental data used to obtain it.

This feedback also closes the gravity problem. The separate distributed
porous-brane return contribution changes the source enclosed at larger radii and
may determine whether a useful local inverse-square region survives, becomes
screened, overshoots, or changes sign. A force law derived while holding an
inconsistent reservoir or return profile fixed is only a provisional local
result.

If the selected branch is relaxational or mixed, the global ledger must also
track entropy, heat, and degraded energy. A permanent particle cannot consume
ordered structure forever while the reservoir and ordered brane remain
unchanged. In a reversible branch, the localized one-way throat drain and the
separate porous-brane return must instead form parts of a coherent globally
closed transport or periodic cycle with complete energy and momentum accounting.
Periodicity of the global cycle does not imply reversal of the throat-mouth flow.

A possible later cosmological extension asks whether the spatially integrated
distributed porous-brane return rate could slowly exceed the integrated
localized throat-drain rate and thereby increase the intrinsic three-volume of
the ordered brane. That idea is downstream of local--global closure. Return does
not by itself imply expansion, accelerated expansion, or a dark-energy analog;
each would require a separate global derivation.

\section{What is postulated and what remains to be shown}

The ontology is intentionally more mature than the derivation status. It defines
a coherent object of study and prevents later calculations from changing the
meaning of the model, but it does not convert open mechanisms into results. The
main status distinctions can be summarized compactly:

\begin{center}
\begin{tabularx}{\textwidth}{@{}L{0.29\textwidth}X@{}}
\toprule
\textbf{Element} & \textbf{Present epistemic status} \\
\midrule
Parent arena and one conserved medium & Postulated: preferred time, four Euclidean spatial dimensions, one compressible substrate, and global constituent conservation. \\
Ordered finite brane and bulk state & The two-state ontology and shear-supporting intent are postulated; a dynamically selected stable slab, microscopic shear origin, and complete spectrum remain open. \\
Light & Two transverse modes and one longitudinal mode follow conditionally from the supplied homogeneous quadratic action; finite-slab guidance, acceptable leakage, and a localized photon packet remain open. \\
Particles & Finite oriented supported throats are the working particle hypothesis; no jointly self-consistent and stable throat solution has yet closed geometry, support, transport, conversion, and reservoir loading. \\
Electric sector & Orientation supplies the postulated sign; the Coulomb-carrying branch, interaction sign, universal coupling, conserved current, complete antiparticle map, and short-range dynamics remain to be derived. \\
Gravity & Localized one-way throat drainage together with separate distributed porous-brane return is the proposed mechanism; inverse-square scaling has conditional precursor support, while the full source reduction, calibration, parity, sign, dynamics, and \(S_{\mathrm{leak}}\) return profile remain open. \\
Magnetism and radiation & Moving-source tensor and falloff structure have conditional precursor support; one-current normalization across electricity, magnetism, and all radiation channels remains a target. \\
Inertia and global closure & Reconfiguration of the complete dressed throat is the proposed origin of inertia; positive response, low-energy universality, acceptable drag, and the local--global reservoir fixed point remain open. \\
\bottomrule
\end{tabularx}
\end{center}

Several claims are deliberately outside this ontology. The model does not claim
to be exact general relativity or exact electromagnetism. It does not insert a
fundamental electromagnetic gauge field merely to recover Maxwell's equations,
and it does not identify the longitudinal brane mode with charge. It does not
claim that support-mode energy is observed mass, that active, passive, and
inertial mass share one microscopic mechanism, or that all characteristic speeds
are automatically equal. It does not yet supply quantum theory, photon number,
particle statistics, or a derivation of the real universe from first principles.

The resulting mental picture is nevertheless unified. One conserved
four-dimensional medium forms an ordered finite-thickness brane inside a
de-structured bulk state. The brane carries transverse light and other collective
modes. Finite oriented throats are candidate particles. Orientation supplies
electric sign, while complete throat structure supplies species. Localized
one-way throat transport, stress, and conversion, together with the separate
distributed porous-brane \(S_{\mathrm{leak}}\) return and reservoir response,
are proposed to underlie gravity. Motion and acceleration produce magnetic,
radiative, and inertial effects. The bulk stores the material and energy
involved and supplies material for distributed reordering into the porous brane;
it does not return material through a throat.

The research program asks whether this entire story can be realized by one
stable parent system without independent retuning. Success requires the same
fields, coefficients, branch choice, source conventions, and reservoir data to
survive every sector. Failure to obtain that common realization would reject the
candidate branch, or ultimately the architecture, even if individual analogies
remain suggestive.
